% SHARD-ID: CA-08-HILBERT-SPACE-COMPILER-SPEC
% SHARD-TITLE: Scoping the Hilbert-Space Compiler
% SHARD-SUMMARY: Specifies the first partial compiler from typed grammar expressions to Hilbert carriers, gradings, observable-algebra policies, symmetry data, and sector witnesses.
% SHARD-SUMMARY: Gives deterministic constructor rules and finite-dimensional invariants that can be checked before any categorical anyon data enters.
% SHARD-SUMMARY: Works through small examples: maybe qubits, distinguishable pairs, derived tensor-power closure, exchange-irrep sectors, invariant quotients, and a required compile error.
% SHARD-KEYWORDS: Hilbert-space compiler, compiler specification, examples, dimensions, observable policy, compile error, symmetric sector, invariant quotient

\section{Scoping the Hilbert-Space Compiler}
\label{sec:hilbert-space-compiler-spec}

\subsection{Status, Sources, and Dependencies}

\textbf{Status: specification shard.}  This shard turns CA-04's contract into a
first partial compiler specification.  Its scope is the pre-work grammar of
CA-02--CA-07: finite Hilbert carriers, derived all-tensor-power completions,
unitary/projective symmetries via central extensions, and irrep/projection sectors.

% Source: literature/md/1910.10619/1910.10619.md:39 -- Fock presentations often mix exchange symmetry and indefinite particle number.
% Source: literature/md/1910.10619/1910.10619.md:41 -- distinguishable Fock first; Bose/Fermi by imposing exchange equivalence.
% Source: literature/md/1910.10619/1910.10619.md:57 -- distinguishable Fock as direct sum of tensor powers.
% Source: literature/md/1910.10619/1910.10619.md:59 -- zero-particle vacuum sector is C.
% Source: references/cft/Schottenloher2008/Schottenloher2008.md:1524 -- unitary representation = continuous homomorphism G -> U(H).
% Source: references/cft/Schottenloher2008/Schottenloher2008.md:1528 -- projective representation = continuous homomorphism G -> U(P).
% Source: references/cft/Schottenloher2008/Schottenloher2008.md:1570 -- projective representation determines a central extension and a lift.
% Source: references/text/PenneysUnitaryFusionCategories.md:219 -- orthogonal direct sums give isometries and mutually orthogonal projections.
% Source: references/text/PenneysUnitaryFusionCategories.md:235 -- orthogonal projections satisfy p^2=p=p^\dagger and split by an isometry.
% Checked: test/runtests.jl, finite C2 invariant projector and two-qubit exchange projectors.

\subsection{Compile Judgement}

The compiler is a \emph{partial} map.  A successful judgement has the form
\[
  \Gamma;\Pi \vdash E \Downarrow
  (H_E,\{H_{E,n}\},\Alg_E,U_E,P_E,Q_E,\mathsf{Ev}_E),
\]
where \(\Gamma\) is the atom table and \(\Pi\) is the policy table.  A failed
judgement has the form
\[
  \Gamma;\Pi \vdash E \Uparrow \mathsf{reason}.
\]
The failure result is part of the specification: incompatible symmetry data,
missing observable policy, or an ill-typed sector constructor must stop compilation.

The atom table entry for \(x\) has at least a Hilbert space \(K_x\), a degree
\(d_x\in\mathbb N\) (usually \(1\)), evidence metadata, and optionally symmetry
data \(U_x:G\to U(K_x)\) or a named projective central extension as in CA-05.

\subsection{Surface Language for the First Compiler}

The first surface language is deliberately small:
\[
\begin{aligned}
E ::= {}& 0 \mid \one \mid x
       \mid \operatorname{OR}(E,F)
       \mid \operatorname{AND}(E,F)
       \mid \operatorname{Sector}_{G,\chi}(E).
\end{aligned}
\]
Here \(G\) acts on the compiled carrier and \(\chi\) is a one-dimensional unitary
character; invariants are the case \(\chi=1\).  Fock, finite truncation,
symmetric sectors, and antisymmetric sectors are derived from OR/AND,
representation propagation, and the generic sector operation.

For a one-particle atom \(x\), define tensor-power expressions recursively by
\[
  x^{[0]}:=\one,\qquad x^{[m+1]}:=\operatorname{AND}(x^{[m]},x).
\]
Then
\[
  \mathsf T_{\le M}(x):=\operatorname{OR}_{m=0}^{M}x^{[m]},\qquad
  \mathsf T_{\infty}(x):=\text{Hilbert completion of } \operatorname{OR}_{m\ge0}x^{[m]}.
\]
\(\mathsf T_{\infty}(x)\) is the usual full/distinguishable Fock space
\(\hoplus_{m\ge0}K_x^{\hotimes m}\), but as a derived graded representation.

\subsection{Constructor Rules}

The carrier, grading, and symmetry rules are recursive:
\[
H_0=\zero,\quad H_{\one}=\mathbb C,\quad H_x=K_x,
\]
\[
H_{\operatorname{OR}(E,F)}=H_E\oplus H_F,\qquad
H_{\operatorname{AND}(E,F)}=H_E\hotimes H_F.
\]
The grading rules are CA-02's:
\[
H_{\operatorname{OR}(E,F),n}=H_{E,n}\oplus H_{F,n},\qquad
H_{\operatorname{AND}(E,F),n}=\bigoplus_{i+j=n}H_{E,i}\hotimes H_{F,j}.
\]
If \(U_E,U_F\) are representations of the same group or named central
extension, the symmetry rules are
\[
U_{\operatorname{OR}}(g)=U_E(g)\oplus U_F(g),\qquad
U_{\operatorname{AND}}(g)=U_E(g)\hotimes U_F(g).
\]
If the symmetry objects are not the same, the carrier may still compile, but the
symmetry output is rejected unless the requested output was explicitly
carrier-only.

For the derived tensor-power closures,
\[
H_{\mathsf T_{\infty}(x)}=\hoplus_{m\ge0}K_x^{\hotimes m},\qquad
H_{\mathsf T_{\le M}(x)}=\bigoplus_{m=0}^{M}K_x^{\otimes m}.
\]
For finite examples the compiler may report dimensions:
\[
\begin{aligned}
\dim\operatorname{OR} &= \dim H_E+\dim H_F,\\
\dim\operatorname{AND} &= (\dim H_E)(\dim H_F),\\
\dim\mathsf T_{\le M}(x) &= \sum_{m=0}^{M}(\dim K_x)^m .
\end{aligned}
\]
These follow from choosing orthonormal bases: disjoint union gives direct-sum
bases, ordered pairs give tensor-product bases, and truncation is a finite sum.

Sector constructors return projection witnesses supplied by representation data.
For finite \(G\), a unitary character \(\chi:G\to U(1)\), and representation
\(U_E:G\to U(H_E)\), set
\[
P_{\chi}=|G|^{-1}\sum_{g\in G}\overline{\chi(g)}\,U_E(g),\qquad
H_{\operatorname{Sector}_{G,\chi}(E)}=\Ran(P_{\chi}).
\]
The quotient relation is generated by \(U_E(g)v-\chi(g)v\).  The invariant
sector is \(\chi=1\).  For exchange sectors on \(K^{\hotimes n}\), \(G=S_n\) and
\(U_E=\Pi_n\); the trivial and sign characters give
\[
P_{{\rm s},n}=n!^{-1}\sum_{\sigma\in S_n}\Pi_n(\sigma),\qquad
P_{{\rm a},n}=n!^{-1}\sum_{\sigma\in S_n}\operatorname{sgn}(\sigma)\Pi_n(\sigma).
\]
The observable algebra is never inferred from these formulas.  It is the policy
slot \(\Pi_{\rm obs}\): carrier-only, full, block, tensor-local, invariant, or compressed.

\subsection{Example 1: A Maybe Qubit}

Let \(Q=\mathbb C^2\) with basis \(e_0,e_1\).  Compile
\[
  E_1=\operatorname{OR}(\one,Q).
\]
The rule gives
\[
  H_{E_1}=\mathbb C\oplus\mathbb C^2,\qquad
  H_{E_1,0}=\mathbb C,\quad H_{E_1,1}=\mathbb C^2,\quad \dim H_{E_1}=3.
\]
With full observable policy, \(\Alg_{E_1}=\BH(\mathbb C\oplus\mathbb C^2)\).
With block policy,
\[
  \Alg_{E_1}^{\rm block}=\BH(\mathbb C)\oplus\BH(\mathbb C^2).
\]
Thus the same carrier has two different compiled observable outputs.

\subsection{Example 2: Two Distinguishable Qubits}

Compile \(E_2=\operatorname{AND}(Q,Q)\).  The carrier is
\[
  H_{E_2}=Q\otimes Q,\qquad
  \{e_0\otimes e_0,e_0\otimes e_1,e_1\otimes e_0,e_1\otimes e_1\}
  \text{ is a basis,}
\]
so \(\dim H_{E_2}=4\), all in degree \(2\).  The tensor-local policy is the
subalgebra generated by
\[
  \BH(Q)\otimes I,\qquad I\otimes \BH(Q),
\]
whereas the full policy is \(\BH(Q\otimes Q)\).  The compiler records which one
was requested; it does not identify them by syntax alone.

\subsection{Example 3: Derived Tensor-Power Closure}

Compile the macro-expanded expression
\[
  E_3=\mathsf T_{\le2}(Q)=\operatorname{OR}(\one,\operatorname{OR}(Q,\operatorname{AND}(Q,Q))).
\]
The output carrier is
\[
  H_{E_3}=\mathbb C\oplus Q\oplus(Q\otimes Q),
\]
with sector dimensions
\[
  \dim H_{0}=1,\qquad \dim H_{1}=2,\qquad \dim H_{2}=4,\qquad
  \dim H_{E_3}=1+2+4=7.
\]
If \(U:G\to U(Q)\) is supplied, symmetry propagation gives
\[
  U_{E_3}(g)=1_{\mathbb C}\oplus U(g)\oplus (U(g)\otimes U(g)).
\]
No Bose/Fermi relation has been imposed yet; this is still the distinguishable
graded representation.

\subsection{Example 4: Two-Qubit Exchange-Irrep Sectors}

Let \(S\) be the swap on \(Q\otimes Q\).  Since
\[
S(e_0e_0)=e_0e_0,\quad S(e_1e_1)=e_1e_1,\quad
S(e_0e_1)=e_1e_0,\quad S(e_1e_0)=e_0e_1,
\]
where \(e_ie_j\) abbreviates \(e_i\otimes e_j\), the exchange representation of
\(S_2\) has two one-dimensional irrep characters, trivial and sign, supplying
projectors with ranges:
\[
  P_{\rm s}=\frac12(I+S),\qquad P_{\rm a}=\frac12(I-S)
\]
\[
  \Ran(P_{\rm s})=\Span\{e_0e_0,\ (e_0e_1+e_1e_0)/\sqrt2,\ e_1e_1\},
\]
\[
  \Ran(P_{\rm a})=\Span\{(e_0e_1-e_1e_0)/\sqrt2\}.
\]
Thus the compiler reports dimensions \(3\) and \(1\) for the trivial- and
sign-irrep sectors.  The Julia test suite checks the projection, orthogonality,
trace/dimension, and basis-vector claims for this example.

\subsection{Example 5: A Finite Invariant Quotient}

Let \(C_2=\{1,s\}\) act on \(H=\mathbb C^2\) by the swap matrix.  Then
\[
  P_{C_2}=\frac12(I+S)
  =\begin{pmatrix}1/2&1/2\\1/2&1/2\end{pmatrix}.
\]
Its range is \(\Span\{(1,1)\}\), and its kernel is
\(\Span\{(1,-1)\}\).  The quotient presentation is therefore
\[
  \mathbb C^2/\Span\{(1,-1)\}\cong \Span\{(1,1)\}.
\]
This is \(\operatorname{Sector}_{C_2,1}(H)\), the trivial-irrep sector of the
\(C_2\) representation.
The invariant observable algebra on the original carrier consists of the
operators commuting with \(S\), while the compressed algebra on the invariant
line is one-dimensional.  These are different observable policies even in this
small example.

\subsection{Example 6: A Required Compile Error}

Let \(A,B\) be atoms with projective symmetries over the same physical group
\(G\), but represented in the input by unrelated central extensions \(E_A\to G\)
and \(E_B\to G\).  The carrier of \(\operatorname{OR}(A,B)\) is still
\(K_A\oplus K_B\).  However, a request for one compiled symmetry representation
must fail:
\[
  \Gamma;\Pi \vdash \operatorname{OR}(A,B)\Uparrow
  \text{``no common symmetry group or central extension specified''.}
\]
This is not a missing feature; it is the strict semantics required by CA-05 and
CONVENTIONS.md (f).

\subsection{Lamport Dependency Skeleton}

{\small
\[
\setlength{\arraycolsep}{2pt}
\begin{array}{rcl}
D24 &:& \Gamma;\Pi\vdash E\Downarrow(H_E,\{H_{E,n}\},\Alg_E,U_E,P_E,Q_E,\mathsf{Ev}_E)
        \text{ or } \Gamma;\Pi\vdash E\Uparrow\mathsf{reason}.\\
D25 &:& E::=0\mid\one\mid x\mid\operatorname{OR}\mid\operatorname{AND}
        \mid\operatorname{Sector}_{G,\chi}.\\
D26 &:& \mathsf T_{\le M}(x)=\operatorname{OR}_{m=0}^{M}x^{[m]},\quad
        H_{\mathsf T_{\le M}}=\bigoplus_{m=0}^{M}K_x^{\otimes m}.\\
D27 &:& P_{\chi}=|G|^{-1}\sum_g\overline{\chi(g)}U_E(g),\quad
        H_{\operatorname{Sector}_{G,\chi}}=\Ran(P_{\chi}).\\
T15 &:& \text{Finite dimensions obey sum, product, and } \sum_{m=0}^{M}d^m
        \text{ by explicit basis constructions.}\\
T16 &:& \text{Irrep/character choices supply projection/subspace/quotient witnesses.}\\
T17 &:& \text{Compiler failure is a specified outcome when requested structure is ill typed or unsupported.}
\end{array}
\]
}

\subsection{Acceptance Criteria and Boundary}

Future implementation work must preserve these acceptance criteria: every
constructor fills typed output slots; every nontrivial mathematical output has an
evidence tag; every quotient names its relation subspace; every observable claim
names its observable policy; and every executable example asserts an invariant,
not merely successful execution.

This compiler does not yet parse fusion-category simples, anyonic braid data,
Hamiltonians, lattice limits, conformal nets, or VOA data.  Those enter only
after their source conventions and invariants are fixed.  CA-08 is the finite
Hilbert-space kernel that later shards must extend without weakening the output
record, the evidence discipline, or the failure semantics.
